% Generated draft: requires mathematical and source review.
\documentclass[11pt]{article}
\usepackage[margin=25mm]{geometry}
\usepackage{fontspec}
\setmainfont{Latin Modern Roman}
\usepackage{amsmath,amssymb,mathtools,unicode-math}
\setmathfont{Latin Modern Math}
\usepackage{enumitem,needspace,xurl,hyperref}
\setlength{\parindent}{0pt}
\setlength{\parskip}{5pt}
\setlength{\emergencystretch}{3em}
\newcommand{\N}{\mathbb{N}}
\newcommand{\Z}{\mathbb{Z}}
\newcommand{\Q}{\mathbb{Q}}
\newcommand{\R}{\mathbb{R}}
\newcommand{\C}{\mathbb{C}}
\newcommand{\F}{\mathbb{F}}
\newcommand{\A}{\mathbb{A}}
\newcommand{\PP}{\mathbb{P}}
\newcommand{\E}{\mathbb{E}}
\newcommand{\eps}{\varepsilon}
\newcommand{\dd}{\,\mathrm{d}}
\begin{document}
\section*{UnsolvedMath ID: 701396}
\textbf{Source problem code:} catalogue entry
\subsection{Definitions and mathematical statement}
Let \(G\) be an arbitrary nontrivial group, let \(Z=\langle t\rangle\cong\Z\) be an infinite cyclic group, and put \(H=Z*G\), the free product of these groups. For \(r\in H\), define its normal closure in \(H\) by
\[
\langle\!\langle r\rangle\!\rangle_H
=\bigl\langle h r h^{-1}:h\in H\bigr\rangle,
\]
the smallest normal subgroup of \(H\) containing \(r\).

The supplied well-known problem asks whether there exist a nontrivial group \(G\) and an element \(r\in Z*G\) such that
\[
\langle\!\langle r\rangle\!\rangle_{Z*G}=Z*G.
\]
Its conjectured negative answer, called Kervaire's conjecture in the supplied editors' comments, is the assertion
\[
\text{for every group }G\ne\{1\}\text{ and every }r\in Z*G,
\qquad
(Z*G)/\langle\!\langle r\rangle\!\rangle_{Z*G}\ne\{1\}.
\]
No finiteness, countability, or torsion hypothesis on \(G\) is imposed in the general question.

The supplied editors' comments report that the negative answer has been proved when \(G\) is torsion-free and when \(G\) is residually finite. Here torsion-free means that \(g^n=1\), for \(g\in G\) and a positive integer \(n\), implies \(g=1\); residually finite means that every nonidentity element of \(G\) has nonidentity image under some homomorphism from \(G\) to a finite group. These special-case results are supplied source qualifications and have not been independently verified here.
\subsection{Short English statement}
Can the free product of an infinite cyclic group and a nontrivial group be generated as a normal subgroup by one element? The supplied conjectured answer is no for every nontrivial group. The editors report this negative answer for torsion-free and residually finite groups.
\subsection{Sources}
\begin{itemize}
\item Local UnsolvedMath registry, permanent ID \(701396\), source code \texttt{catalogue entry}, title \emph{Kourovka Notebook Problem 17.94}. The supplied registry labels the problem open; this label is not an independently verified literature assessment.
\item \emph{The Kourovka Notebook (Unsolved Problems in Group Theory)}, original problem 17.94, 17th issue (2010); supplied source edition: 21st issue (2026), \texttt{arXiv:1401.0300v48}. The collection identifies E. I. Khukhro and V. D. Mazurov as editors. Proposed by J. Mycielski in 2010; the statement describes it as a well-known problem.
\item Supplied editors' reference for the torsion-free case: A. Klyachko, \emph{Commun. Algebra} \textbf{21}, no. 7 (1993), 2555--2575.
\item Supplied editors' reference for the residually finite case: M. Gerstenhaber and O. S. Rothaus, \emph{Proc. Nat. Acad. Sci. U.S.A.} \textbf{48} (1962), 1531--1533.
\end{itemize}
\end{document}
